\subsection{Partial basis and the \texorpdfstring{$J_2$}{J2} example}
\label{sec:partial-basis-j2}

The following example distinguishes a count of individual coordinate
intertwiners from the dimension of the full intertwining space. It also
shows why the action on the constant quotient need not be scalar.
Let $L=\mathbb F_q(\theta)$ and put
\[
 J_2=\begin{pmatrix}0&1\\0&0\end{pmatrix},\qquad
 \Phi_t=\begin{pmatrix}\theta+\tau^2&1\\0&\theta+\tau^2\end{pmatrix},
 \qquad
 \Psi_t=\begin{pmatrix}\theta&1+\tau\\0&\theta\end{pmatrix}.
\]
Thus $N_\Phi=N_\Psi=J_2$. The leading coefficient of $\Phi_t$ is
$I_2$, so we may take $C=I_2$, $\sigma=(1,2)$ and
$(j_1,j_2)=(2,2)$. The degree gap holds:
$\deg_\tau\Psi_t=1<2$.
The full normal representatives have degree less than two in every entry.
Consequently, $M=8$, with ordered basis
\begin{equation}
 \mathcal E=(E_{11},E_{11}\tau,E_{12},E_{12}\tau,
             E_{21},E_{21}\tau,E_{22},E_{22}\tau).
 \label{eq:partial-j2-full-basis}
\end{equation}
The extra $\tau$ in $\Psi_t$ acts on a scalar coefficient $c$
by $c\mapsto c^q$. Terms of degree two are reduced using inner
biderivations. In the order \eqref{eq:partial-j2-full-basis}, the full action is
\begin{equation}
 \Pi_t=\begin{pmatrix}\theta I_4&A(\tau)\\0&\theta I_4\end{pmatrix},
 \qquad
 A(\tau)=\begin{pmatrix}
  1&0&0&0\\
  \tau&1&0&0\\
  0&-\tau&1&0\\
  0&0&\tau&1
 \end{pmatrix}.
 \label{eq:partial-j2-full-action}
\end{equation}
For example, the degree-two term in $\Psi_t(cE_{21}\tau)$
is $c^qE_{11}\tau^2$, whose normal reduction is $-c^qE_{12}$.
This accounts for the entry $-\tau$ in $A(\tau)$.

\paragraph{The constant map and the individual direction count.}
For $H=\begin{pmatrix}x&y\\z&w\end{pmatrix}$, the constant map is
\begin{equation}
 \mathcal T(H)=HJ_2-J_2H
     =\begin{pmatrix}-z&x-w\\0&z\end{pmatrix}.
 \label{eq:partial-j2-constant-map}
\end{equation}
With column vectorization $\operatorname{vec}(H)=(x,z,y,w)^T$, its matrix is
\[
 D=\begin{pmatrix}
  0&-1&0&0\\
  0&0&0&0\\
  1&0&0&-1\\
  0&1&0&0
 \end{pmatrix}.
\]
Hence
\begin{equation}
 \operatorname{rank}D=2,\qquad
 \ker\mathcal T=L I_2\oplus L E_{12},\qquad
 h=4-\operatorname{rank}D=2.
 \label{eq:partial-j2-rank}
\end{equation}
Among the four individual matrix units, only $E_{12}$ lies in this
kernel. The second kernel direction, $I_2=E_{11}+E_{22}$, is a linear
combination of two units whose images are nonzero. Thus the rule
$\varepsilon_{i,s}=0$ when $E_{i,s}J_2=J_2E_{i,s}$, and $1$ otherwise,
gives
\[
 \varepsilon=\begin{pmatrix}1&0\\1&1\end{pmatrix}.
\]
The number of zero entries is one, whereas the full kernel dimension
is two. A dimension calculation using the individual count would predict
$8-1=7$; the correct dimension of
$\operatorname{Ext}_{0,\tau}(\Phi,\Psi)$ is $8-h=6$.

This discrepancy is visible directly in the candidate partial basis.
The shifted degree ranges $1\leq r<j_s+\varepsilon_{i,s}$ give
\[
 (E_{11}\tau,E_{11}\tau^2,E_{12}\tau,
  E_{21}\tau,E_{21}\tau^2,E_{22}\tau,E_{22}\tau^2).
\]
These seven representatives do not give seven independent additive
coordinates. For a constant matrix $H$,
\[
 \delta^{(H)}_t=H\Phi_t-\Psi_tH
   =\mathcal T(H)+H\tau^2-J_2H^{(1)}\tau,
\]
where $H^{(1)}$ denotes entrywise $q$-th powers. In particular, for every
$c\in L$,
\[
 \delta^{(cI_2)}_t=cI_2\tau^2-c^qE_{12}\tau.
\]
Consequently,
\begin{equation}
 [cE_{11}\tau^2]+[cE_{22}\tau^2]-[c^qE_{12}\tau]=0.
 \label{eq:partial-j2-dependence}
\end{equation}
In the displayed candidate order, the additive parameter vector
$(0,c,-c^q,0,0,0,c)$ therefore represents zero. The missing intertwiner
$I_2$ produces this one-parameter redundancy, missed by counting individual
coordinate directions. The family parametrizes the subgroup but is not a
coordinate basis; the partial basis applicability test rejects it.

\paragraph{The true subgroup in normal coordinates.}
The image of the constant map is
\[
 \operatorname{im}\mathcal T
 =L E_{12}\oplus L(E_{22}-E_{11}).
\]
The rank-based construction consequently gives the ordered basis
\begin{equation}
 \mathcal E_0=(E_{12},E_{22}-E_{11},
               E_{11}\tau,E_{12}\tau,E_{21}\tau,E_{22}\tau).
 \label{eq:partial-j2-zero-basis}
\end{equation}
The first two normal representatives have nonzero constant terms, but
these constants are removable by inner biderivations. Indeed,
\[
 [cE_{12}]=[-cE_{11}\tau^2],\qquad
 [c(E_{22}-E_{11})]=[c^qE_{11}\tau-cE_{21}\tau^2]
 \quad(c\in L).
\]
Both classes therefore admit representatives with zero constant term,
as required by the definition of $\operatorname{Ext}_{0,\tau}$.
In the order \eqref{eq:partial-j2-zero-basis}, the restricted action is
\begin{equation}
 \Pi_{0,t}=\begin{pmatrix}
  \theta&1&0&0&-\tau&0\\
  0&\theta&0&0&0&0\\
  0&0&\theta&0&1&0\\
  0&\tau&0&\theta&0&1\\
  0&0&0&0&\theta&0\\
  0&0&0&0&0&\theta
 \end{pmatrix}.
 \label{eq:partial-j2-zero-action}
\end{equation}
The inclusion, expressed in the full normal coordinates, is the constant map
\begin{equation}
 K(u_1,\ldots,u_6)^T
   =(-u_2,u_3,u_1,u_4,0,u_5,u_2,u_6)^T.
 \label{eq:partial-j2-inclusion}
\end{equation}

\paragraph{The nonscalar quotient and the exact sequence.}
Choose $[E_{21}]$ and $[E_{22}]$ as the ordered basis of
$\operatorname{Mat}_2(L)/\operatorname{im}\mathcal T$.
For $X=\begin{pmatrix}x&y\\z&w\end{pmatrix}$, its quotient coordinates
are $(z,x+w)^T$. In the full coordinate order
\eqref{eq:partial-j2-full-basis}, the projection is therefore
\begin{equation}
 P(x_1,\ldots,x_8)^T=(x_5,x_1+x_7)^T.
 \label{eq:partial-j2-projection}
\end{equation}
The induced action of $J_2$ sends
$[E_{21}]$ to $[E_{11}]=[E_{22}]$ and $[E_{22}]$ to
$[E_{12}]=0$. Thus the quotient action is
\begin{equation}
 Q_t=\begin{pmatrix}\theta&0\\1&\theta\end{pmatrix},
 \qquad Q_t\neq\theta I_2.
 \label{eq:partial-j2-quotient}
\end{equation}
The maps above satisfy
\[
 \Pi_tK=K\Pi_{0,t},\qquad P\Pi_t=Q_tP,\qquad PK=0.
\]
Moreover, $K$ is injective, $P$ is surjective and $\ker P=\operatorname{im}K$.
These facts follow immediately from
\eqref{eq:partial-j2-inclusion}--\eqref{eq:partial-j2-projection}.
These identities are understood in the Ore ring $L\{\tau\}$.
The entries of $K$ and $P$ are $0,1,-1$, hence Frobenius-fixed, so the
identities can be checked coefficientwise despite the $\tau$-terms in the
action matrices.
We obtain the exact sequence of $t$-modules
\begin{equation}
 0\longrightarrow(\mathbb G_a^6,\Pi_{0,t})
 \xrightarrow{K}(\mathbb G_a^8,\Pi_t)
 \xrightarrow{P}(\mathbb G_a^2,Q_t)
 \longrightarrow0.
 \label{eq:partial-j2-exact-sequence}
\end{equation}
In terms of extension classes, its first two terms are
$\operatorname{Ext}_{0,\tau}(\Phi,\Psi)$ and
$\operatorname{Ext}^1_\tau(\Phi,\Psi)$, respectively.
The differentials of $\Pi_t$, $\Pi_{0,t}$ and $Q_t$ are $\theta I$
plus square-zero nilpotent matrices.
The computations and the dependence
\eqref{eq:partial-j2-dependence} remain valid in every characteristic,
including characteristic two.

This example shows that the full rank calculation and the induced quotient
action are essential: the number of individual coordinate intertwiners
need not equal $h$, and the constant quotient need not have scalar
$t$-action. The supplied \texttt{Counterexample-J2.nb} computes this pair,
its bases, the exact-sequence maps and the partial basis applicability
conditions.
